\section{De la corrección de código a las tareas de seguridad: transferencia de capacidades y reconfiguración de escenarios}

\subsection{Diferencias estructurales entre las tareas de seguridad y las tareas de corrección comunes}
Las tareas de seguridad no se limitan a «hacer que las pruebas pasen»; también exigen responder tres preguntas más difíciles: si la vulnerabilidad es realmente activable, cuál es la condición que la activa y si puede reproducirse en distintas rutas. En comparación con la corrección de código en general, la búsqueda de vulnerabilidades depende más de la observación del comportamiento dinámico, del razonamiento sobre el flujo de datos entre funciones y de la verificación contrafactual.

\begin{knowledgebox}{El ciclo mínimo de tareas de un Agent de seguridad}
\begin{enumerate}
    \item Suponer que cierta ruta podría tener una condición de desbordamiento, acceso fuera de límites o falta de verificación;
    \item generar una entrada ejecutable o un script que la active;
    \item ejecutarla bajo un sanitizer o depurador y recopilar evidencia;
    \item si falla, corregir la suposición y seguir explorando hasta reproducir o refutar el caso.
\end{enumerate}
\end{knowledgebox}

\subsection{CTF/Enigma como escalón de transición}
La clase presenta las tareas de tipo CTF y Enigma como una capa intermedia: enfatizan más las cadenas de explotación y activación que los parches cotidianos dentro del IDE, pero siguen siendo más controlables que las vulnerabilidades de sistemas reales a gran escala. Este escalón tiene valor porque permite al equipo iterar sus herramientas y su flujo de control en un entorno evaluable.

\begin{importantbox}{Por qué empezar con CTF antes de pasar a vulnerabilidades reales}
El CTF permite verificar rápidamente, en un entorno reproducible, si «la herramienta sirve»; las vulnerabilidades reales, en cambio, están más influidas por el ruido del contexto, el entorno de dependencias y la complejidad del espacio de entradas. Pulir primero el ciclo básico del Agent y luego migrar hacia sistemas reales produce una tasa de éxito más alta.
\end{importantbox}

\subsection{Las herramientas de análisis dinámico entran al flujo principal}
En la segunda mitad del minuto 00:50, el presentador habló específicamente del papel del sanitizer como analizador dinámico. A diferencia de un juicio puramente estático, el sanitizer ofrece directamente, en tiempo de ejecución, la señal de «si se activó un comportamiento de memoria ilegal», lo que lo hace naturalmente apto para ser el oracle del Agent.

\begin{warningbox}{Depender solo de un clasificador estático lleva a caer en «parece una vulnerabilidad»}
Muchas funciones parecen peligrosas de forma local, pero la cadena de llamadas global puede ya incluir verificaciones de límites; a la inversa, algo que parece seguro localmente puede seguir siendo explotable al combinarse entre módulos. Sin verificación dinámica, tanto los falsos positivos como los falsos negativos aumentan de forma considerable.
\end{warningbox}

\subsection{Resumen de la sección}
Las tareas de seguridad no son una simple ampliación de las tareas de corrección de código, sino un cambio de la función objetivo. Buscar solo que las pruebas pasen no puede responder si «realmente existe una vulnerabilidad explotable»; es necesario introducir evidencia de ejecución dinámica y un ciclo de verificación más sólido.

